\documentclass[10pt,a4paper]{amsart}
\usepackage[margin=19mm]{geometry}
\usepackage{amsmath,amssymb,mathtools}
\usepackage[scr=rsfs,cal=euler]{mathalfa}
\usepackage{xfrac}
\usepackage[unicode,hidelinks,bookmarks=false]{hyperref}
\newcommand{\ie}{i.e.\ }
\newcommand{\cf}{cf.\ }
\newcommand{\resp}{respectively\ }
\newcommand{\Subsection}{\subsection}
\providecommand{\nobreakdash}{\nobreak\hbox{-}}
\setlength{\emergencystretch}{2em}
\multlinegap0pt
\usepackage{xcolor}
\usepackage{mathtools}
\usepackage{amssymb}
\usepackage[shortlabels]{enumitem}
\usepackage{tikz}
\usepackage{tensor}
\newtheorem{definition}{Definition}
\newtheorem{lemma}{Lemma}
\newtheorem{proposition}{Proposition}
\usepackage{cleveref}
\crefname{definition}{Definition}{Definitions}
\crefname{lemma}{Lemma}{Lemmas}
\crefname{proposition}{Proposition}{Propositions}
\newcommand{\G}{\mathcal{G}}
\title{Quantum metrics from length functions on \'etale groupoids}
\author{Are Austad}
\date{Source: arXiv:2602.20032v2 (CC BY 4.0)}
\begin{document}
\maketitle

\noindent\emph{Source and licence.} Quantum metrics from length functions on \'etale groupoids. Version arXiv:2602.20032v2 (CC BY 4.0), distributed by arXiv under the Creative Commons Attribution 4.0 International licence, http://creativecommons.org/licenses/by/4.0/, as recorded in the arXiv licence metadata for this exact version and shown on the abstract page https://arxiv.org/abs/2602.20032v2. Full licence notice: \texttt{licenses/arxiv-2602.20032-CC-BY-4.0.txt}.\par\smallskip

\noindent\textit{Editorial scope.} Counted unit: the article's proposition prop:ellB-is-cont-length-function (source0.tex lines 1550--1555). The article's complete original proof of it is delivered with it (source0.tex lines 1557--1573, one \texttt{proof} environment of 17 lines). No mathematics is written, repaired or re-derived: the statement and the proof are the article's own lines, in the article's own notation, and no retained line is substituted or re-flowed.  Retained uncounted as the source-local context the proof needs: lines 1489--1491; lines 1513--1526; lines 1520--1525; lines 1532--1536.  Nothing was omitted from any retained span, so the package carries no omission record.  No retained line contains a citation, so the wrapper carries no bibliography and no bibliography entry.  No retained line contains a figure environment, an image-inclusion command or drawing code, so no image is delivered and the package declares no asset dependency.  Editorial formatting only, in addition: an AMS \texttt{amsart} preamble that restates the article's own declarations and macros used by the retained lines, namely the environments \texttt{definition}, \texttt{lemma}, \texttt{proposition} on separate counters, together with the standard packages those lines need.  The retained environments print in this document as Definition 1, Lemma 1, Proposition 1 in order of appearance, whereas the article prints the same items with the numbering of its own full document.  The complete native source of the article as distributed in the arXiv e-print for this exact version is bundled unchanged as \texttt{sources/195153/source0.tex}.
	\begin{align}\label{eq:def-cylinder-set}
		Z(\mu, \lambda) = \{ (x,y) \in Z(\mu) \times Z(\lambda) \mid x_{[m + \vert \mu \vert + 1, \infty)} = y_{[n + \vert \lambda \vert + 1 , \infty)} \}.
	\end{align}

	\begin{definition}\label{def:AF-length-function} 
	Let $\G$ be an AF groupoid arising from a Bratteli diagram $B = (V,E)$ with $\vert S(B)\vert < \infty$. 
	For $(x,y) \in \G$, denote by 
	\begin{align*}
		k(x,y) = \min \{m \mid x_{[m+1,\infty)} = y_{[m+1,\infty)}\}.
	\end{align*}
	 Moreover, we define the function
	\begin{align}\label{eq:def-AF-length-function}
		\ell(x,y) = \begin{cases}
			0 & \text{ if  $x=y$} \\ 
			\sum_{v \in S(B)} \vert P(v, V_{k(x,y)}) \vert & \text{ if $x \neq y$}.
		\end{cases}
	\end{align}
\end{definition}

	\begin{align}
		\ell(x,y) = \begin{cases}
			0 & \text{ if  $x=y$} \\ 
			\sum_{v \in S(B)} \vert P(v, V_{k(x,y)}) \vert & \text{ if $x \neq y$}.
		\end{cases}
	\end{align}

	\begin{lemma}\label{lemma:relating-ell-and-k}
		Let $\G$ be an AF groupoid arising from a Bratteli diagram $B = (V,E)$ with $\vert S(B)\vert < \infty$. 
		Let $k$ and $\ell$ be as in \cref{def:AF-length-function}. 
		Suppose $(x,y), (z,w) \in \G$ with $x\neq y$ and $z\neq w$. If $\ell(x,y) > \ell(z,w)$, then $k(x,y) > k(z,w)$. 
	\end{lemma}

	\begin{proposition}\label{prop:ellB-is-cont-length-function}
		Let $\G$ be an AF groupoid arising from a Bratteli diagram $B = (V,E)$ with $\vert S(B)\vert < \infty$.  
		The function $\ell$ from \eqref{eq:def-AF-length-function} 
		is a continuous length function on 
		$\G$.  
	\end{proposition}

	\begin{proof}
		We first verify that $\ell$ is a length function. 
		First, note that $\ell(x,x) = 0$ by definition. Moreover, if $x\neq y$, then $k(x,y)\geq 1$ and so $P(v, V_{k(x,y)}) \neq \emptyset$ for at least one $v \in S(B)$. We conclude that $\ell(x,y) = 0$ if and only if $x = y$. 
		Since $(x,y)^{-1} = (y,x)$ in $\G$, we see that $\ell((x,y)^{-1}) =  \ell(x,y)$, since $k(x,y) = k(y,x)$. 
		Now, let $(x,y), (y,z), (x,z) \in \G$. We wish to show that
		\begin{align*}
			\ell(x,z) \leq \ell(x,y) + \ell(y,z).
		\end{align*}
		Suppose for the sake of contradiction that $\ell(x,z) > \max \{ \ell(x,y), \ell(y,z)\}$. By \cref{lemma:relating-ell-and-k} we then have $k(x,z) > \max \{ k(x,y), k(y,z) \}$. But then $x_{[l_1, \infty)} = y_{[l_1, \infty)}$ for some $l_1 < k(x,z)$, and $y_{[l_2, \infty)} = z_{[l_2, \infty)}$ for some $l_2 < k(x,z)$.   We conclude that $x_{[l, \infty)} = z_{[l, \infty)}$ for $l < k(x,z)$ contradicting the minimality of $k(x,z)$. We deduce a stronger version of the triangle inequality, namely $\ell(x,z) \leq \max \{ \ell(x,y) , \ell(y,z) \}$. 
 
		
		To see that $\ell$ is continuous, suppose $(x^{(j)}, y^{(j)}) \to (x,y)$ in $\G$. There are finite paths $\mu, \lambda$ with $t(\mu) = t(\lambda)$ such that $(x,y) \in Z(\mu, \lambda)$, cf. \eqref{eq:def-cylinder-set}. We then have $x_{[m +\vert \mu  \vert + 1 , \infty)} = y_{[n + \vert \lambda \vert + 1 , \infty)}$, and suppose that $m + \vert \mu \vert$ is minimal with this property, that is, $x$ and $y$ do not start agreeing at an earlier level. Since $Z(\mu, \lambda)$ is open, there is $N$ such that for all $j \geq N$, $(x^{(j)}, y^{(j)}) \in Z(\mu, \lambda)$. We claim that $N$ can also be chosen so that $x^{(j)}$ and $y^{(j)}$ do not agree at a level earlier than $m + \vert \mu \vert = n + \vert \lambda \vert$. Indeed, if no such $N$ exists, that is, if for arbitrarily large $j$, there is $l_j < m + \vert \mu \vert +1$ such that $x^{(j)}_{[l_j, \infty)} =  y^{(j)}_{[l_j, \infty)}$, we would have
		\begin{align*}
			s(x^{(j)}, y^{(j)}) \not \to s(x,y) = y \quad \text{or} \quad r(x^{(j)}, y^{(j)}) \not \to r(x,y) = x,
		\end{align*}
		thus implying that $s$ or $r$ is not continuous, a contradiction. Thus $x^{(j)}$ and $y^{(j)}$ eventually agree from level $m + \vert \mu \vert $, and $\ell (x^{(j)}, y^{(j)}) = \ell (x,y)$. Thus $\ell$ is continuous. 
	\end{proof}

\end{document}
